% !TEX root = ../../thesis.tex
\chapter{Deep Learning Architectures for Semantic Segmentation}
\label{cap:architectures}

Semantic segmentation in agricultural robotics consists in assigning a predefined category label to every pixel of an acquired image. For autonomous row navigation, the perception system must distinguish navigable soil corridors from surrounding crop canopies, weeds, and field boundaries under variable outdoor illumination.

This chapter presents the mathematical formulation and architectural design of the deep learning models evaluated in this work. First, the segmentation task and encoder-decoder principles are defined. Second, Fully Convolutional Networks (FCN) and DeepLabV3 are detailed with emphasis on atrous convolutions and multi-scale context aggregation. Third, MobileNetV3 and the Lite Reduced Atrous Spatial Pyramid Pooling (LR-ASPP) head are analyzed, detailing the depthwise separable convolutions and channel attention mechanisms that reduce computational complexity for edge computing. Finally, the Segment Anything Model 3 (SAM~3) is described as an offline annotation engine for training label generation.

\section{Problem formulation and encoder-decoder principles}
\label{sec:seg_formulation}

Let $I \in \mathbb{R}^{H \times W \times 3}$ denote an input color image with height $H$ and width $W$. Semantic segmentation aims to estimate a dense label map $Y \in \{1, \dots, C\}^{H \times W}$, where each pixel $(i,j)$ is assigned to one of $C$ mutually exclusive semantic classes. Deep neural networks approximate this mapping through a parameterized function $f_\theta(I)$ that outputs an unnormalized score tensor $Z \in \mathbb{R}^{H \times W \times C}$. The posterior probability of class $c$ at pixel $(i,j)$ is obtained via the spatial softmax function:
\begin{equation}
    p_{i,j,c} = \frac{\exp(Z_{i,j,c})}{\sum_{k=1}^{C} \exp(Z_{i,j,k})}.
    \label{eq:softmax}
\end{equation}

During supervised training with ground-truth one-hot masks $Y^* \in \{0,1\}^{H \times W \times C}$, parameters $\theta$ are optimized by minimizing the multi-class cross-entropy loss:
\begin{equation}
    \mathcal{L}_{\text{CE}}(Y^*, p) = -\frac{1}{HW} \sum_{i=1}^{H} \sum_{j=1}^{W} \sum_{c=1}^{C} Y^*_{i,j,c} \log(p_{i,j,c}).
    \label{eq:ce_loss}
\end{equation}

Standard convolutional architectures for image classification progressively reduce spatial dimensions via strided convolutions and max-pooling layers. While spatial reduction increases the receptive field and aggregates contextual information, it discards fine-grained boundary positions. Semantic segmentation architectures solve this trade-off using an encoder-decoder topology:
\begin{itemize}
    \item \textbf{Encoder (Backbone):} A sequence of convolutional stages that extracts high-level semantic features $F_{\text{enc}} \in \mathbb{R}^{\frac{H}{S} \times \frac{W}{S} \times C_{\text{enc}}}$, where $S \ge 1$ denotes the total output stride (typically $S \in \{8, 16, 32\}$).
    \item \textbf{Decoder (Segmentation Head):} A module that restores spatial resolution to $H \times W$ through bilinear interpolation or transposed convolutions, combining low-level spatial details from earlier encoder layers with abstract high-level features.
\end{itemize}

\section{Fully Convolutional Networks (FCN)}
\label{sec:fcn_arch}

Long et al.~\cite{long2015fcn} introduced Fully Convolutional Networks by replacing the final dense (fully connected) classification layers of standard convolutional backbones with convolutional filters.

\subsection{Fully convolutional transformation}
\label{subsec:fcn_conv}

In standard classification networks, the spatial feature map produced by the final convolutional layer is flattened into a vector and passed to dense layers, restricting the network to fixed input resolutions. FCN replaces each dense layer with a $1 \times 1$ convolutional layer whose kernel depth matches the dense weight matrix. For a feature map $F \in \mathbb{R}^{h \times w \times C_{\text{in}}}$, a convolutional layer with $C$ filters of size $1 \times 1$ computes:
\begin{equation}
    Z_{i,j,c} = b_c + \sum_{k=1}^{C_{\text{in}}} W_{c,k} \cdot F_{i,j,k}, \qquad \forall i \in \{1,\dots,h\}, \; j \in \{1,\dots,w\},
    \label{eq:conv1x1}
\end{equation}
producing a coarse score map of dimension $h \times w \times C$. This allows the network to process inputs of arbitrary dimensions in a single forward pass without sliding-window recomputation.

\subsection{Skip architecture and upsampling}
\label{subsec:fcn_skip}

To restore the coarse score map from stride $S=32$ back to the original image dimensions $H \times W$, FCN employs transposed convolutions (also known as backward strided convolutions or fractionally strided convolutions). A transposed convolution with stride $s$ places the input values on a dilated grid and computes the output via standard convolution, learning an upsampling filter $w \in \mathbb{R}^{K \times K}$:
\begin{equation}
    \hat{y}_{i',j'} = \sum_{u=0}^{K-1} \sum_{v=0}^{K-1} w_{u,v} \cdot x_{\lfloor i'/s \rfloor - u, \; \lfloor j'/s \rfloor - v}.
    \label{eq:transposed_conv}
\end{equation}
Upsampling directly from stride 32 produces blurry boundaries. Long et al. therefore designed skip connections:
\begin{itemize}
    \item \textbf{FCN-32s:} Upsamples the final stride-32 score map directly by a factor of 32 in a single step.
    \item \textbf{FCN-16s:} Combines the $2\times$ upsampled stride-32 score map with predictions extracted from the stride-16 stage (e.g., \texttt{conv4}), then applies $16\times$ upsampling.
    \item \textbf{FCN-8s:} Combines the $2\times$ upsampled FCN-16s score map with predictions extracted from the stride-8 stage (e.g., \texttt{conv3}), followed by an $8\times$ upsampling step.
\end{itemize}
The implementation used in this work couples an FCN segmentation head with a ResNet-50 backbone.

\section{DeepLabV3 and atrous convolutions}
\label{sec:deeplabv3_arch}

DeepLabV3, developed by Chen et al.~\cite{chen2017deeplab}, addresses the resolution loss caused by consecutive pooling operations by replacing standard downsampling with atrous (dilated) convolutions and multi-scale pooling modules.

\subsection{Atrous convolution formulation}
\label{subsec:atrous_conv}

Let $x[i,j]$ denote a 2D input feature map and $w[p,q]$ a convolutional filter of size $K \times K$ with $p, q \in \{-M, \dots, M\}$ where $M = \frac{K-1}{2}$. Standard convolution samples adjacent positions on the input grid. Atrous convolution introduces a dilation rate $r \in \mathbb{N}^+$, defining the step with which the filter samples the input signal:
\begin{equation}
    y[i, j] = \sum_{p=-M}^{M} \sum_{q=-M}^{M} x[i + r \cdot p, \; j + r \cdot q] \cdot w[p, q].
    \label{eq:atrous_conv}
\end{equation}

When $r = 1$, Eq.~\eqref{eq:atrous_conv} reduces to the standard convolution. For $r > 1$, the filter samples inputs with spaces of $r - 1$ zeros between consecutive filter taps. The effective filter size $K'$ becomes:
\begin{equation}
    K' = K + (K - 1)(r - 1).
    \label{eq:effective_kernel}
\end{equation}
For example, a $3 \times 3$ kernel ($K=3$) with dilation rate $r = 6$ covers a spatial footprint of $3 + 2 \times 5 = 13 \times 13$ pixels, requiring only 9 learnable parameters. Atrous convolution expands the receptive field without adding parameters or reducing the spatial resolution of the feature maps.

\subsection{Atrous Spatial Pyramid Pooling (ASPP)}
\label{subsec:aspp_module}

Objects and agricultural corridors appear at varying visual scales depending on camera mounting angle and distance from the vehicle. DeepLabV3 captures contextual information across multiple scales using the Atrous Spatial Pyramid Pooling (ASPP) module.

The ASPP module operates in parallel on the high-level feature map extracted by the backbone (typically at output stride $S=16$) and consists of five branches:
\begin{enumerate}
    \item One $1 \times 1$ standard convolution with 256 filters and batch normalization.
    \item Three $3 \times 3$ atrous convolutions with 256 filters each, using dilation rates $r \in \{6, 12, 18\}$ for an output stride of 16 (or $r \in \{12, 24, 36\}$ for stride 8).
    \item One Image-Level Pooling branch: to incorporate global image context, this branch computes global average pooling across the spatial dimensions of the input tensor, processes the resulting $1 \times 1 \times C_{\text{in}}$ vector with a $1 \times 1$ convolution (256 filters) and batch normalization, and upsamples the output bilinearly to the original ASPP spatial resolution.
\end{enumerate}

The outputs of all five branches are concatenated along the channel dimension, resulting in a tensor with $5 \times 256 = 1280$ channels:
\begin{equation}
    F_{\text{concat}} = \left[ F_{1\times1}, \; F_{\text{atrous}, r_1}, \; F_{\text{atrous}, r_2}, \; F_{\text{atrous}, r_3}, \; F_{\text{image\_pool}} \right].
    \label{eq:aspp_concat}
\end{equation}
Finally, a $1 \times 1$ convolution with 256 filters, batch normalization, and dropout (with probability 0.5) projects $F_{\text{concat}}$ to 256 channels, followed by a $1 \times 1$ classification layer and bilinear upsampling to image resolution $H \times W$.

\section{MobileNetV3 and the Lite Reduced ASPP (LR-ASPP) head}
\label{sec:mobilenet_lraspp}

While DeepLabV3 with a ResNet-50 backbone provides high segmentation accuracy, its 50.8 million parameters and heavy computational requirement result in high inference latency (86.87~ms per image), exceeding the real-time budget of mobile robot platforms. To achieve real-time operation on edge hardware, Howard et al.~\cite{howard2019mobilenetv3} developed MobileNetV3 and the Lite Reduced ASPP (LR-ASPP) segmentation head.

\subsection{Depthwise separable convolutions}
\label{subsec:dw_conv}

The fundamental building block of MobileNet architectures is the depthwise separable convolution, which factors a standard convolution into two separate steps: a spatial depthwise convolution and a channelwise pointwise convolution.

Consider an input feature map $X \in \mathbb{R}^{H \times W \times C_{\text{in}}}$ and an output feature map $Y \in \mathbb{R}^{H \times W \times C_{\text{out}}}$, processed with spatial filter size $D_K \times D_K$.
\begin{itemize}
    \item \textbf{Standard convolution:} Each filter channel interacts with every input channel simultaneously. The computational cost, measured in Multiply-Accumulate (MAC) operations, is:
    \begin{equation}
        \text{Cost}_{\text{standard}} = H \cdot W \cdot C_{\text{in}} \cdot C_{\text{out}} \cdot D_K^2.
        \label{eq:cost_standard}
    \end{equation}
    \item \textbf{Depthwise convolution:} Applies a single $D_K \times D_K$ spatial filter per input channel independently, without channel cross-talk:
    \begin{equation}
        \text{Cost}_{\text{depthwise}} = H \cdot W \cdot C_{\text{in}} \cdot D_K^2.
        \label{eq:cost_dw}
    \end{equation}
    \item \textbf{Pointwise convolution:} Applies $C_{\text{out}}$ filters of size $1 \times 1$ across all channels to compute linear combinations of the depthwise outputs:
    \begin{equation}
        \text{Cost}_{\text{pointwise}} = H \cdot W \cdot C_{\text{in}} \cdot C_{\text{out}}.
        \label{eq:cost_pw}
    \end{equation}
\end{itemize}
The total cost of a depthwise separable convolution is the sum of Eqs.~\eqref{eq:cost_dw} and~\eqref{eq:cost_pw}:
\begin{equation}
    \text{Cost}_{\text{sep}} = H \cdot W \cdot C_{\text{in}} \cdot D_K^2 + H \cdot W \cdot C_{\text{in}} \cdot C_{\text{out}}.
    \label{eq:cost_sep}
\end{equation}
Comparing Eq.~\eqref{eq:cost_sep} to Eq.~\eqref{eq:cost_standard} yields the computational reduction ratio:
\begin{equation}
    \frac{\text{Cost}_{\text{sep}}}{\text{Cost}_{\text{standard}}} = \frac{H \cdot W \cdot C_{\text{in}} \cdot D_K^2 + H \cdot W \cdot C_{\text{in}} \cdot C_{\text{out}}}{H \cdot W \cdot C_{\text{in}} \cdot C_{\text{out}} \cdot D_K^2} = \frac{1}{C_{\text{out}}} + \frac{1}{D_K^2}.
    \label{eq:reduction_ratio}
\end{equation}
For standard $3 \times 3$ filters ($D_K = 3$) and typical channel counts ($C_{\text{out}} \ge 64$), the reduction ratio is approximately $\frac{1}{9} \approx 0.11$. Depthwise separable convolutions thus reduce arithmetic operations and parameter counts by roughly 88--89\% relative to standard convolutions.

\subsection{Inverted residual blocks with linear bottlenecks}
\label{subsec:inverted_residuals}

In classical residual blocks (such as those in ResNet), the input channels are first compressed via a $1 \times 1$ convolution, processed with a $3 \times 3$ convolution, and expanded back to the original channel count (narrow $\to$ wide $\to$ narrow).

MobileNetV3 reverses this structure using inverted residual blocks (narrow $\to$ wide $\to$ narrow):
\begin{enumerate}
    \item \textbf{Expansion layer:} A $1 \times 1$ pointwise convolution expands the low-dimensional input $C_{\text{in}}$ to a higher-dimensional space $e \cdot C_{\text{in}}$, where $e \in \{3, 6\}$ is the expansion factor.
    \item \textbf{Depthwise layer:} A $3 \times 3$ or $5 \times 5$ depthwise separable convolution performs spatial feature extraction in the expanded channel space.
    \item \textbf{Linear bottleneck:} A $1 \times 1$ projection layer compresses the channels back to $C_{\text{out}}$. Crucially, no non-linear activation function is placed after this final projection layer. Because non-linearities such as ReLU discard negative values by projecting them to zero, applying non-linear activations in low-dimensional spaces destroys useful geometric information. Retaining a linear bottleneck preserves manifold topology.
\end{enumerate}
When the input and output channel dimensions match and the stride is 1, a residual shortcut connection adds the input directly to the block output.

\subsection{Squeeze-and-Excitation (SE) channel attention}
\label{subsec:se_attention}

To adaptively calibrate channel dependencies, MobileNetV3 incorporates Squeeze-and-Excitation (SE) blocks within the inverted residual structure:
\begin{enumerate}
    \item \textbf{Squeeze:} Global spatial average pooling aggregates the spatial dimensions of feature map $U \in \mathbb{R}^{H' \times W' \times C'}$ into a channel descriptor vector $z \in \mathbb{R}^{C'}$:
    \begin{equation}
        z_c = \frac{1}{H' W'} \sum_{i=1}^{H'} \sum_{j=1}^{W'} U_{i,j,c}.
        \label{eq:se_squeeze}
    \end{equation}
    \item \textbf{Excitation:} Two fully connected layers parameterized by weights $W_1 \in \mathbb{R}^{\frac{C'}{4} \times C'}$ and $W_2 \in \mathbb{R}^{C' \times \frac{C'}{4}}$ model non-linear channel relationships with a channel reduction ratio of 4:
    \begin{equation}
        s = \sigma\left(W_2 \cdot \text{ReLU}(W_1 z)\right),
        \label{eq:se_excite}
    \end{equation}
    where $\sigma(\cdot)$ denotes the hard-sigmoid activation function.
    \item \textbf{Scale:} The input tensor $U$ is rescaled channel-wise by the activation vector $s$:
    \begin{equation}
        \tilde{U}_{i,j,c} = s_c \cdot U_{i,j,c}, \qquad \forall c \in \{1, \dots, C'\}.
        \label{eq:se_scale}
    \end{equation}
\end{enumerate}

\subsection{Efficient non-linearities: hard-swish and hard-sigmoid}
\label{subsec:hard_activations}

The swish non-linearity, defined as $\text{swish}(x) = x \cdot \sigma(x)$, improves gradient flow in deep networks compared to standard ReLU. However, computing the exponential function $\sigma(x) = (1 + e^{-x})^{-1}$ introduces noticeable latency on embedded mobile processors lacking dedicated transcendental hardware units.

MobileNetV3 replaces standard sigmoid and swish with piece-wise linear hard approximations based on $\text{ReLU6}(x) = \min(\max(x, 0), 6)$:
\begin{equation}
    \text{h-sigmoid}(x) = \frac{\text{ReLU6}(x + 3)}{6},
    \label{eq:h_sigmoid}
\end{equation}
\begin{equation}
    \text{h-swish}(x) = x \cdot \text{h-sigmoid}(x) = x \cdot \frac{\text{ReLU6}(x + 3)}{6}.
    \label{eq:h_swish}
\end{equation}
Eqs.~\eqref{eq:h_sigmoid} and~\eqref{eq:h_swish} eliminate floating-point exponential operations, preserving the performance benefits of swish while minimizing execution time on embedded CPUs and GPUs.

\subsection{The Lite Reduced ASPP (LR-ASPP) architecture}
\label{subsec:lraspp_arch}

Standard DeepLabV3 ASPP requires multiple parallel atrous branches with 256 channels each, resulting in significant memory bandwidth requirements. Howard et al. proposed the Lite Reduced ASPP (LR-ASPP) head as a lightweight alternative tailored for mobile segmentation.

The LR-ASPP architecture uses features from two distinct stages of the MobileNetV3-Large backbone:
\begin{itemize}
    \item A low-resolution, high-semantic feature map $F_{\text{low}} \in \mathbb{R}^{\frac{H}{16} \times \frac{W}{16} \times 960}$ extracted at output stride 16.
    \item A high-resolution, low-semantic feature map $F_{\text{high}} \in \mathbb{R}^{\frac{H}{8} \times \frac{W}{8} \times 40}$ extracted at output stride 8.
\end{itemize}

The processing steps of LR-ASPP proceed as follows:
\begin{enumerate}
    \item $F_{\text{low}}$ is projected to 128 channels using a $1 \times 1$ convolution with batch normalization:
    \begin{equation}
        V_{\text{low}} = \text{Conv}_{1\times1}(F_{\text{low}}) \in \mathbb{R}^{\frac{H}{16} \times \frac{W}{16} \times 128}.
    \end{equation}
    \item A global context branch computes global average pooling across $F_{\text{low}}$ using a large pooling kernel (spatial size matching the feature map), followed by a $1 \times 1$ convolution to 128 channels and hard-sigmoid activation:
    \begin{equation}
        G = \text{h-sigmoid}\left(\text{Conv}_{1\times1}(\text{AvgPool}(F_{\text{low}}))\right) \in \mathbb{R}^{1 \times 1 \times 128}.
    \end{equation}
    \item The low-resolution feature map $V_{\text{low}}$ is modulated by the global context gating vector $G$ via element-wise multiplication:
    \begin{equation}
        M_{\text{low}} = V_{\text{low}} \odot G.
    \end{equation}
    \item The modulated tensor $M_{\text{low}}$ is upsampled by a factor of 2 via bilinear interpolation to match the stride-8 resolution of $F_{\text{high}}$.
    \item $F_{\text{high}}$ is projected to 128 channels via a $1 \times 1$ convolution and added element-wise to the upsampled low-resolution features:
    \begin{equation}
        F_{\text{fused}} = \text{BilinearUpsample}_{2\times}(M_{\text{low}}) + \text{Conv}_{1\times1}(F_{\text{high}}).
        \label{eq:lraspp_fusion}
    \end{equation}
    \item A final $1 \times 1$ convolutional classifier projects $F_{\text{fused}}$ to $C$ class logits, followed by an $8\times$ bilinear upsampling step back to the original image dimensions $H \times W$.
\end{enumerate}

By discarding multiple parallel dilated convolutions in favor of a single gated global pooling path and direct stride-8 feature addition, LR-ASPP reduces computational operations while retaining global scene context and boundary localization.

\section{Segment Anything Model 3 (SAM~3) for dataset annotation}
\label{sec:sam3_theory}

Manually annotating agricultural images at pixel precision requires considerable labor: drawing polygons around irregular plant boundaries, inter-row soil, and weed patches requires between 15 and 30 minutes per frame. Foundation models provide a mechanism to scale up annotation through zero-shot visual prompting.

Kirillov et al.~\cite{kirillov2023sam} established the promptable segmentation paradigm with SAM, using Vision Transformers to generate binary segmentation masks from foreground/background points, bounding boxes, or coarse free-form masks. Carion et al.~\cite{carion2026sam3} extended this architecture in SAM~3 by integrating open-vocabulary semantic concepts and text-prompted segmentation.

\subsection{Architecture overview}
\label{subsec:sam3_arch}

SAM~3 consists of three primary components:
\begin{enumerate}
    \item \textbf{Image Encoder:} A heavy Vision Transformer (ViT-H or ViT-L) pre-trained with masked autoencoding and contrastive objectives. The image encoder processes an input image $I \in \mathbb{R}^{1024 \times 1024 \times 3}$ and outputs a downsampled $64 \times 64 \times 256$ spatial image embedding. Computing this embedding represents the primary computational cost, but it is executed once per image regardless of the number of subsequent prompts.
    \item \textbf{Prompt Encoder:} Embeds sparse prompts (spatial coordinate points, bounding boxes) through positional encodings summed with learned prompt-type embeddings. Text prompts are embedded using a frozen contrastive text encoder (such as CLIP or an aligned concept encoder) that maps target category descriptions (e.g., \texttt{"lettuce plant"}, \texttt{"navigable soil path"}) to the prompt feature space.
    \item \textbf{Two-Way Mask Decoder:} A lightweight transformer decoder that alternates self-attention among prompt tokens and cross-attention between prompt tokens and the image embedding in both directions. The updated tokens are mapped via linear layers to dynamic filter kernels, which are applied to an upsampled feature map via point-wise dot product to produce candidate segmentation masks and associated IoU confidence scores.
\end{enumerate}

\subsection{Operational distinction from onboard models}
\label{subsec:sam3_role}

A key architectural distinction must be emphasized between the foundation model and the robotic segmentation network:
\begin{itemize}
    \item SAM~3 requires over 600 million parameters and multi-gigabyte memory footprints, with single-image forward times exceeding 1000~ms on workstation GPUs. It cannot be deployed directly for closed-loop visual servoing on resource-constrained embedded computers such as the Jetson platform.
    \item In this thesis, SAM~3 functions strictly as an \textbf{offline supervision generator}. Given text prompts representing crop and corridor concepts, SAM~3 generates dense candidate masks across lettuce, tomato, and cucumber image collections.
    \item The resulting masks are reviewed, filtered, and used as ground-truth targets to train lightweight real-time architectures (FCN, DeepLabV3, and LR-ASPP). The distilled lightweight models execute inference on the robot platform in 10 to 86~ms, maintaining real-time control capability.
\end{itemize}

\section{Summary of architectural characteristics}
\label{sec:arch_summary}

Table~\ref{tab:architecture_characteristics} summarizes the theoretical characteristics of the four architectures analyzed in this chapter, establishing the foundation for the dataset experiments and comparative evaluations presented in Chapter~\ref{cap:dataset_training}.

\begin{table}[htbp]
    \centering
    \small
    \setlength{\tabcolsep}{4pt}
    \caption{Theoretical comparison of deep learning architectures evaluated in this thesis.}
    \label{tab:architecture_characteristics}
    \begin{tabularx}{\textwidth}{@{}l >{\raggedright\arraybackslash}p{2.6cm} >{\raggedright\arraybackslash}X >{\raggedright\arraybackslash}p{2.6cm} >{\raggedright\arraybackslash}X@{}}
        \toprule
        Model & Backbone & Context mechanism & Target role & Primary constraint \\
        \midrule
        FCN & ResNet-50 & Transposed convolutions and skip connections & Benchmark segmentation & High latency (58.24~ms), coarse boundaries \\
        \addlinespace
        DeepLabV3 & ResNet-50 & Parallel Atrous Spatial Pyramid Pooling (ASPP) & High-accuracy benchmark & Highest latency (86.87~ms), 50.8M parameters \\
        \addlinespace
        LR-ASPP & MobileNetV3-Large & Gated global average pooling and stride-8 fusion & Real-time robotic navigation & Compact capacity, fine texture handling \\
        \addlinespace
        SAM~3 & ViT-Large & Foundation promptable Vision Transformer & Offline automatic annotation & Non-real-time ($>1$~s), high GPU memory footprint \\
        \bottomrule
    \end{tabularx}
\end{table}
